\documentclass[11pt]{article}

% Change "review" to "final" to generate the final (sometimes called camera-ready) version.
% Change to "preprint" to generate a non-anonymous version with page numbers.
\usepackage[preprint]{acl}

\usepackage{float}
\usepackage{dblfloatfix}



% Standard package includes
\usepackage{times}
\usepackage{latexsym}

% For proper rendering and hyphenation of words containing Latin characters (including in bib files)
\usepackage[T1]{fontenc}
% This assumes your files are encoded as UTF8
\usepackage[utf8]{inputenc}

% This is not strictly necessary, and may be commented out,
% but it will improve the layout of the manuscript,
% and will typically save some space.
\usepackage{microtype}

% This is also not strictly necessary, and may be commented out.
% However, it will improve the aesthetics of text in
% the typewriter font.
\usepackage{inconsolata}

% Including images in your LaTeX document requires adding
% additional package(s)
\usepackage{graphicx}

% If the title and author information does not fit in the area allocated, uncomment the following
%
%\setlength\titlebox{<dim>}
%
% and set <dim> to something 5cm or larger.

\title{Fine-Tuning Pre-Trained Transformers for Systematic Compositional Generalization on SCAN Tasks}

% Author information - update with your details
\author{Mats Ellingsen \\
  University of Copenhagen \\
  Department of Computer Science \\
  \texttt{srh673@alumni.ku.dk}}

\begin{document}
\maketitle
\begin{abstract}
This report investigates how a fine-tuned pre-trained language model compares to a Transformer trained from scratch under the same training conditions. Following Lake and Baroni's (2018) SCAN experiments, I compare a fine-tuned GPT-2 model against our group's encoder-decoder Transformer baseline on the simple split (Experiment 1b) and length split (Experiment 2). Results show that GPT-2 consistently underperforms the baseline, with a ~30 percentage-point gap on in-distribution generalization and substantially worse token accuracy on length extrapolation. These findings suggest that pre-training on natural language does not transfer to the rule-based compositional reasoning required by SCAN tasks, indicating that natural language composition and algebraic compositionality operate through fundamentally different computational mechanisms.
\end{abstract}

\section{Introduction}

The ability to understand and produce novel combinations from known components, I.e. systematic compositionality, is a hallmark of human language and cognition \citep{Fodor:88}. \citet{lake2018generalizationsystematicitycompositionalskills} demonstrated that while recurrent neural networks (RNNs) excel at generalizing to random subsets of compositional commands through pattern matching, they fail catastrophically when systematic compositional reasoning is required. Their SCAN (Simplified CommAI Navigation) tasks provide a controlled testbed for evaluating whether neural networks can achieve algebraic compositionality.

Modern Transformers have achieved remarkable success in language understanding and generation tasks. However, our group's reimplementation of Lake and Baroni's experiments using encoder-decoder Transformers revealed similar systematicity failures as the original RNN models, particularly on length extrapolation and primitive command generalization tasks.

This raises a critical question: how do pre-trained language models, having been exposed to diverse compositional patterns in natural language, compare to models trained from scratch on systematic compositional generalization? This report investigates this question by comparing a fine-tuned GPT-2 model against our group's encoder-decoder Transformer baseline under the same training conditions on SCAN tasks.

\section{Background and Related Work}

\subsection{Systematicity in Neural Networks}

The systematicity debate, initiated by \citet{Fodor:88}, questions whether neural networks can capture the compositional structure of human cognition. \citet{lake2018generalizationsystematicitycompositionalskills} operationalized this debate through the SCAN domain, where models must map natural language commands (e.g., ``jump twice and walk around left'') to action sequences (e.g., JUMP JUMP LTURN WALK LTURN WALK).

Their experiments revealed that RNNs succeed when test examples involve ``mix-and-match'' recombination of training patterns but fail when systematic rule application is required. For instance, models that observed all compositional modifiers with primitives like ``walk'' and ``run'' could not generalize them to ``jump'' when the latter was only seen in isolation during training.

\subsection{SCAN Tasks and Splits}

SCAN pairs brief navigation commands with action sequences, enabling controlled train--test splits that probe different facets of compositionality:

\textbf{Experiment 1 (Simple Split).} Random train--test split where most compositional patterns appear during training, measuring generalization under broad distribution overlap with varying fractions of training data (1\%--100\%) to probe data efficiency.

\textbf{Experiment 2 (Length Split).} Train on commands whose target action sequences are $\leq 22$ steps; test on sequences of 24--48 steps, probing systematic length extrapolation. This experiment has two evaluation conditions: (1) \textit{without oracle}, where the model must autonomously determine output length, and (2) \textit{with oracle}, where a length constraint forces the model to generate exactly the target sequence length. The oracle condition isolates compositional structure learning from length prediction, revealing whether models can compose learned patterns over longer horizons when length ambiguity is removed. sequence-level exact match and token-level accuracy is reported for both conditions.

\textbf{Experiment 3 (Add-Primitive Split).} Hold out a primitive command from all composed examples during training while still showing it in isolation, testing whether compositional modifiers generalize to this held-out primitive under varying numbers of composed examples.

\subsection{Group Reimplementation Results}

We reimplemented the three SCAN experiments with an encoder--decoder Transformer in PyTorch, following the original splits. For Experiment 1, separate models were trained from scratch for each training data fraction. These results serve as the baseline for comparison with the fine-tuned GPT-2 model evaluated in this report.

\textbf{Experiment 1 (Simple Split).} Sequence accuracy scales with the fraction of training commands: $0.63$ at $1\%$, $0.81$ at $2\%$, $0.93$ at $4\%$, and plateaus around $0.97$ once $\geq16\%$ of commands are seen. The model reaches near-ceiling performance with substantially less data than the full corpus.

\textit{Analysis:} This near-perfect performance demonstrates that Transformers excel at in-distribution generalization when compositional patterns overlap between training and test sets. However, this success is driven primarily by recombination of statistically similar components rather than systematic compositional rules. The rapid learning curve---achieving near-perfect accuracy with only 16\% of training data---suggests the model efficiently captures surface-level patterns through attention mechanisms. This mirrors the original findings of \citet{lake2018generalizationsystematicitycompositionalskills} that neural networks succeed on ``mix-and-match'' recombination tasks where novel combinations remain statistically similar to training examples.

\textbf{Experiment 2 (Length Split).} Without oracle-length decoding, sequence accuracy collapses (0 for source lengths 4--7 and effectively 0 for target lengths $\geq26$), despite token accuracy remaining $\sim0.63$--$0.77$. With oracle length, short commands reach perfect sequence accuracy (1.0 at source length 4) but extrapolation remains poor: sequence accuracy drops from $0.85$ at target length 24 to $0.13$ at 48, even though token accuracy stays $\geq0.72$. The gap between sequence- and token-level scores indicates the model learns local patterns but fails to compose them over longer horizons.

\textit{Analysis:} The sharp performance drop for longer sequences reveals a critical limitation: the model exhibits limited ability to systematically scale learned rules beyond training distribution. While token-level accuracy remains moderate ($\sim$0.63--0.77), the near-zero sequence accuracy without oracle indicates complete failure to predict appropriate output length. Even with oracle-length decoding---which removes length prediction uncertainty---performance degrades severely from 0.85 at length 24 to 0.13 at length 48. This dramatic sequence-token gap demonstrates that the model learns local token patterns and short-range dependencies but cannot recursively compose them into correct longer sequences, indicating reliance on statistical regularities rather than systematic compositional rules.

\textbf{Experiment 3 (Add-Primitive Split).} When ``jump'' is only seen in isolation, sequence accuracy is $\approx0.002$ (token $0.64$), while ``turn left'' is better but remains low (sequence $0.29$, token $0.79$). Adding more composed examples raises performance but does not close the gap: with 32 composed instances, sequence accuracy is $0.46$ (token $0.95$), and even with extensive exposure (99 examples) sequence accuracy stays around $0.32$. This mirrors the original finding that Transformers struggle to apply modifiers to unseen primitives.

\textit{Analysis:} The severe failure to generalize to held-out primitives reveals that the model is mostly pattern-matching rather than truly learning compositional rules. The extremely low sequence accuracy for ``jump'' (0.13\%) despite moderate token accuracy (63.5\%) demonstrates that the model does not learn modifiers as algebraic functions that transfer to new primitives. While ``turn left'' shows better performance (sequence accuracy $\sim$30\%, token accuracy $\sim$79\%), this asymmetry reflects differences in surface-level learnability rather than systematic compositionality. Critically, as more composed examples are added to training, token-level accuracy improves steadily while sequence-level accuracy remains very low until many examples are seen. This pattern indicates that improvement is driven by exposure to specific primitive-modifier combinations rather than by learning abstract compositional rules, providing strong evidence that generalization relies on distributional similarity rather than systematic rule application.

\textbf{Summary of Group Results:} Our reimplementation successfully reproduces Lake and Baroni's key findings: Transformers achieve excellent in-distribution generalization (Experiment 1) but fail catastrophically when systematic compositionality is required for length extrapolation (Experiment 2) or primitive generalization (Experiment 3). These results establish the baseline against which we evaluate whether pre-training provides advantages for systematic compositional generalization. In Section 5, these baseline results are presented side-by-side with the fine-tuned GPT-2 results to facilitate direct comparison.

\section{Definition of the Project}

\subsection{Project Objective}

This project evaluates whether GPT-2's pre-training on diverse natural language provides advantages for systematic compositional generalization on the simple split (Experiment 1b) and length split (Experiment 2), comparing against the group's encoder-decoder Transformer baseline under controlled experimental conditions.

\subsection{Approach and Methodology}

I employ GPT-2 (small), a pre-trained decoder-only Transformer, because its autoregressive architecture aligns with SCAN's conditional sequence generation. The SCAN task is reformulated as conditional sequence generation by concatenating commands and actions (\texttt{command <SEP> actions <EOS>}), enabling GPT-2 to generate action sequences autoregressively conditioned on the command prefix while reusing its pre-training on diverse natural language. Training and evaluation mirror the group's encoder-decoder baseline to isolate the effect of pre-training (see Section~\ref{sec:implementation} for exact settings).

\subsection{Implementation}
\label{sec:implementation}
GPT-2 base (124M parameters) was fine-tuned using the HuggingFace Transformers library. For Experiment 1b, eight separate models were trained---one for each training data fraction (1\%, 2\%, 4\%, 8\%, 16\%, 32\%, 64\%, 100\%)---with each model initialized from pre-trained GPT-2 weights. Training used batch size 64 with 1,562 steps (simple split) and batch size 16 with 6,250 steps (length split), with a 15-epoch cap. Hyperparameters: AdamW with learning rates 7e-4 (simple split) and 2e-4 (length split), 500 warmup steps, weight decay 0.01, gradient clipping 1.0. Evaluation used greedy decoding with batch size 128. For Experiment 2's oracle condition, a custom \texttt{LogitsProcessor} enforced target sequence length. Due to compute constraints, single-run evaluation is reported rather than averaging over multiple runs.

\section{Results}

\subsection{Experiment 1b: Simple Split}

Figure~\ref{fig:exp1b_comparison} presents the token accuracy as a function of training data fraction, comparing the encoder-decoder Transformer baseline against the fine-tuned GPT-2 model. Both models follow a similar upward trend as training data increases, but with markedly different scales. The baseline exhibits rapid saturation: starting at 63\% (1\% data), it climbs steeply to 81\% (2\%), 93\% (4\%), and 96\% (8\%), then plateaus near 97\% for all larger fractions (16\% onwards). GPT-2 also improves monotonically but at a much slower rate and lower absolute level: 32\% (1\%) → 36\% (2\%) → 47\% (4\%) → 58\% (8\%) → 65\% (16\%) → 66-67\% (32-100\%), showing continued modest gains beyond 8\% data before flattening around 67\%. The baseline maintains a consistent advantage of approximately 30 percentage points across all fractions.

\begin{figure}[t]
  \centering
  \includegraphics[width=0.7\columnwidth]{figures/exp1b_comparison.png}
  \caption{Token accuracy on Experiment 1b (Simple Split) across different training data fractions. Grouped bars show encoder-decoder Transformer baseline (left bar) versus fine-tuned GPT-2 (right bar) for each training percentage.}
  \label{fig:exp1b_comparison}
\end{figure}


\subsection{Experiment 2: Length Split}

Figure~\ref{fig:exp2_no_oracle} presents results without oracle-length decoding across all metrics: sequence and token accuracy for both source and target lengths. Figure~\ref{fig:exp2_with_oracle} shows the corresponding results with oracle-length decoding.

\textbf{Without Oracle (Figure~\ref{fig:exp2_no_oracle}):} Both models exhibit near-zero sequence accuracy across all source and target lengths. For source lengths (4--9), the baseline achieves 0\% for lengths 4--7, then rises marginally to 7.6\% (length 8) and 4.1\% (length 9), while GPT-2 remains at 0\% across all lengths. For target lengths (24--48), the baseline reaches 33.5\% at length 24, then drops to near-zero (0--3.4\%) for all longer lengths, while GPT-2 achieves 0\% uniformly across all target lengths. Token accuracy shows: the baseline maintains moderate performance (62--91\% across target lengths), while GPT-2 ranges from 17--33\% for target lengths.

\textbf{With Oracle (Figure~\ref{fig:exp2_with_oracle}):} For sequence accuracy by source length, the baseline shows a peak (100\% at length 4) followed by a plateau around 45--57\% for longer source lengths (6--9), while GPT-2 remains near-zero for lengths 4--7 and rises only to 9.4\% (length 8) and 10\% (length 9). By target length, the baseline starts at 84.7\% (length 24), gradually declining to 47--59\% for mid-range lengths (25--32), then dropping to near-zero for very long sequences (36--48). GPT-2 follows a descending trend at lower values: 42.9\% (length 24), 37.3\% (length 25), then 0\% for all longer targets. Token accuracy by target length shows the baseline between 77--97\% while GPT-2 ranges from 42--92\%.

\FloatBarrier

\begin{figure*}[t]
  \centering
  \includegraphics[width=0.7\textwidth]{figures/exp2_no_oracle_all.png}
  \caption{Experiment 2 (Length Split) without oracle-length decoding. Top row: source length results (left: sequence accuracy, right: token accuracy). Bottom row: target length results (left: sequence accuracy, right: token accuracy). Grouped bars show encoder-decoder Transformer baseline versus fine-tuned GPT-2.}
  \label{fig:exp2_no_oracle}
\end{figure*}

\begin{figure*}[t]
  \centering
  \includegraphics[width=0.7\textwidth]{figures/exp2_with_oracle_all.png}
  \caption{Experiment 2 (Length Split) with oracle-length decoding. Top row: source length results (left: sequence accuracy, right: token accuracy). Bottom row: target length results (left: sequence accuracy, right: token accuracy). Grouped bars show encoder-decoder Transformer baseline versus fine-tuned GPT-2.}
  \label{fig:exp2_with_oracle}
\end{figure*}



\section{Discussion}

The central question of this investigation was whether pre-training on natural language corpora provides transferable compositional knowledge for systematic generalization on SCAN tasks. The results indicate that fine-tuned GPT-2 consistently underperforms the encoder-decoder Transformer baseline across both experiments, suggesting pre-training does not confer advantages for the algebraic compositionality required by SCAN.

\subsection{Interpretation of Results}

The 30 percentage-point gap between baseline and GPT-2 across all training fractions in Experiment 1b reveals a fundamental mismatch. While both models exhibit similar upward trends with increased data, the baseline's rapid saturation at 97\% demonstrates superior pattern-matching efficiency. In Experiment 2, both models fail catastrophically at sequence-level generalization without oracle-length guidance (near-zero sequence accuracy). However, the baseline's moderate token accuracy (62--91\%) versus GPT-2's poor performance (17--33\%) indicates the baseline captures more robust local patterns. With oracle-length decoding, the parallel decline in sequence accuracy as target length increases suggests shared architectural limitations, yet the baseline retains partial sequence accuracy at mid-range lengths (47--59\% for lengths 25--32) while GPT-2 collapses to 0\% beyond length 25.

\subsection{Methodological Confounds}

Two critical methodological factors likely explain GPT-2's underperformance, independent of theoretical questions about pre-training's utility.

\textbf{Fine-Tuning on Small Datasets.} Pre-trained language models are optimized for large task-specific datasets that enable gradual adaptation. SCAN's training sets, particularly at low data fractions (1\%--16\%), contain only hundreds to thousands of examples---orders of magnitude smaller than typical fine-tuning scenarios. Fine-tuning methods may struggle to reorient GPT-2's 124M parameters trained on billions of natural language tokens toward SCAN's narrow domain. The baseline, trained from scratch, faces no such parameter reorientation challenge. This dataset size mismatch represents a critical confound: GPT-2's underperformance may reflect fine-tuning methodology limitations rather than fundamental inadequacy of pre-trained representations.

\textbf{Suboptimal Hyperparameters.} The hyperparameters used for GPT-2 fine-tuning were adapted from the baseline's training configuration rather than optimized for fine-tuning. Pre-trained models typically require lower learning rates or techniques like gradual unfreezing to prevent catastrophic forgetting. The lack of systematic hyperparameter search---constrained by compute limitations---means the reported results may substantially underestimate GPT-2's potential.

Beyond methodology, architectural differences contribute to the gap. The encoder-decoder baseline's bidirectional encoder can attend to entire commands simultaneously, while GPT-2's decoder-only architecture processes concatenated sequences autoregressively, treating commands and actions uniformly rather than modeling their conditional relationship explicitly.

\subsection{Implications for Systematic Compositionality}

Both models exhibit the sequence-token accuracy gap characteristic of pattern-matching over systematic rules: they learn local token transitions but fail to compose these patterns into correct longer sequences. This shared failure---even with oracle-length decoding---suggests that neither Transformer variant captures algebraic compositionality as defined by \citet{Fodor:88}. The consistent 30-point gap in Experiment 1b and parallel decline curves in Experiment 2 indicate pre-training does not fundamentally alter the systematicity profile. The absence of advantages on low-data fractions or better length extrapolation suggests natural language pre-training and SCAN compositionality tap into distinct computational mechanisms: SCAN requires algebraic rule application, whereas natural language models learn statistical approximations that generalize through distributional similarity.

\section{Conclusion}

This investigation compared fine-tuned GPT-2 against an encoder-decoder Transformer baseline on SCAN tasks. GPT-2 consistently underperforms with a ~30 percentage-point gap on in-distribution generalization (Experiment 1b) and substantially worse token accuracy on length extrapolation (Experiment 2). Both models exhibit pattern-matching behavior, learning local transitions but failing to compose them into longer sequences, even with oracle-length decoding.

These results suggest that pre-training on natural language does not transfer to the algebraic compositionality required by SCAN, likely because natural language composition---driven by statistical regularities---operates through fundamentally different mechanisms than SCAN's deterministic rules. Achieving human-like systematic compositionality might require architectural innovations beyond standard Transformers rather than relying on pre-training alone.

% Bibliography - create a custom.bib file with your references
\bibliography{custom}

\end{document}
